\chapter{Opportunity Identification}
\label{section:opportunity identification part 1}
% this is some introduction about the topic
\subsection{Pre-class Activities}

As the name suggests, this is an activity that each member should do prior to the class meeting.
\subsubsection{Activity 1: Generate Charter}
\textbf{Objective:} \\
The objective of this activity is to guide you in creating and refining an innovation charter for your team project.

The innovation charter serves as a document that articulates the goals and boundary conditions for your innovation effort.
For example, the charter for the FroliCat effort was:


\begin{quote}
    "Create a physical product in the cat toy category that we can launch to the market within about a year through our existing retail sales channel."
    
    \hfill --charter for the FroliCat
\end{quote}

To begin the task,

\begin{enumerate}
    \item Draft one or more innovation charters that you believe will be advantageous for your project.
    \begin{enumerate}
        \item Keep in mind that it is recommended to create a charter with a broad scope which will allow for a broad focus in the innovation charter and subsequently opens up opportunities that the team may not have considered initially.
        \item Also, ensure the charter to be broader than what the team may initially feel comfortable with as this will challenge your team's assumptions about the types of opportunities that your team should pursue.
    \end{enumerate}
    \item After drafting, it is advisable to review and refine them as necessary.
    \begin{enumerate}
        \item Pay attention to the clarity of the goals and the boundary conditions established in the charters.
    \end{enumerate}
    \item Assume that there are several innovation chart, choose the most suitable innovation charter.
    \begin{enumerate}
        \item Consider the goals and boundary conditions that align best with your project's objectives and constraints.
    \end{enumerate}
\end{enumerate}

\subsection{In-class Activities}
\subsubsection{Activity 1:  Generate Charter}
During this stage, it is expected that each group member has individually generated several innovation charters according to the given guidelines.
\begin{enumerate}
    \item Discuss with all members whether it is more advantageous to select the best innovation charter that had been drafted by one of the members, or to combine elements from multiple charters into a concise sentence.
    \begin{enumerate}
        \item Consider the diverse perspectives of each group member during the discussion.
        \item Evaluate the benefits and drawbacks of choosing a single charter versus combining elements from multiple charters by take into account the shared vision, goals, and desired outcomes of the project.
    \end{enumerate}
    \item Once consensus is reached, draft a new innovation charter that reflects the collective agreement of all team members.
    \begin{enumerate}
        \item Ensure that the new innovation charter is clear, concise, and representative of the team vision and objectives.
        \item Review and refine the drafted innovation charter to ensure that it aligns with the desired outcomes and sets appropriate boundaries for the innovation effort.
    \end{enumerate}
\end{enumerate}